% GENERATED by build_chapters.py; edit docs/report/chapters, not this file

\chapter{GitHub Actions and the Notify Job}

Everything on the GitHub side is one workflow file, \texttt{.\allowbreak{}github/\allowbreak{}workflows/\allowbreak{}demo-\allowbreak{}ci.\allowbreak{}yml}, named "Demo CI", plus the small Node application it tests in \texttt{demo-\allowbreak{}app/\allowbreak{}}. The workflow does two jobs: it tests the code, and when the tests fail because of a code bug it wakes the healer. This chapter goes through the file from top to bottom.

\section{Triggers}

\begin{lstlisting}
on:
  push:
    branches:
      - main
    paths: ["demo-app/**"]
  pull_request:
    branches:
      - main
    paths: ["demo-app/**"]
\end{lstlisting}

(source: .github/\allowbreak{}workflows/\allowbreak{}demo-ci.yml)

The workflow starts in two cases. A push to \texttt{main} starts it, and so does a pull request that targets \texttt{main}. In both cases the changed files must include something under \texttt{demo-\allowbreak{}app/\allowbreak{}}. The path filter keeps documentation commits, report files and the rest of the repository from running the tests.

The path filter has a side effect. A change that touches only files outside \texttt{demo-\allowbreak{}app/\allowbreak{}} does not start CI at all, so a failing test could not be reported for it. That does not matter for the demo, whose bugs live in the app, but it would matter for a real repository, where the filter would have to cover every folder that can break the build.

The \texttt{pull\_\allowbreak{}request} trigger was added so that the fix pull requests opened by the healer are tested too. The reason is in the last section of this chapter.

\section{Jobs}

There are two jobs, \texttt{build-\allowbreak{}and-\allowbreak{}test} and \texttt{notify-\allowbreak{}n8n-\allowbreak{}on-\allowbreak{}failure}. The second one names the first in \texttt{needs}, so it waits for it and can read its outputs. Both run on \texttt{ubuntu-\allowbreak{}latest}.

\section{The build-and-test job}

The job's default working directory is \texttt{demo-\allowbreak{}app}. Its steps are as follows (source: .github/\allowbreak{}workflows/\allowbreak{}demo-ci.yml).

\begin{enumerate}[leftmargin=*, itemsep=2pt]
\item Check out the code with \texttt{actions/\allowbreak{}checkout@v4}.
\item Install Node.js 20 with \texttt{actions/\allowbreak{}setup-\allowbreak{}node@v4}.
\item Run \texttt{npm install}.
\item Run the tests (step id \texttt{tests}) and classify the result.
\item Run \texttt{npm run build -\allowbreak{}-\allowbreak{}if-\allowbreak{}present}.
\end{enumerate}

The fourth step is the one that matters for healing. It runs \texttt{npm test} with \texttt{set +e}, so the script survives a failure and can inspect the exit code:

\begin{lstlisting}
- name: Run tests
  id: tests
  run: |
    set +e
    npm test
    code=$?
    set -e
    if [ "$code" -eq 0 ]; then
      kind=none
    elif [ "$code" -eq 75 ]; then
      kind=infra
    else
      kind=code
    fi
    echo "failure_kind=$kind" >> "$GITHUB_OUTPUT"
    exit "$code"
\end{lstlisting}

(shortened; source: .github/\allowbreak{}workflows/\allowbreak{}demo-ci.yml, step "Run tests")

The result is published as a job output, \texttt{failure\_\allowbreak{}kind}, which the second job reads. The step then exits with the original code, so the job still turns red when the tests fail. The exit code 75 is \texttt{EX\_\allowbreak{}TEMPFAIL}, the conventional "try again later" code (source: demo-app/\allowbreak{}test.js, header comment).

Two situations leave \texttt{failure\_\allowbreak{}kind} without the value \texttt{code}, and both are consequences of how the step is written, not of anything the workflow says. If \texttt{npm install} fails, the tests step never runs, so the output is empty. If the tests pass and only the build step fails, the output is \texttt{none}. Either way the notify job's condition is false and the healer is not called. The first case is intended, since a failed install is an environment matter. The second is a gap: a build failure after passing tests is not healed. In this demo the build script only runs \texttt{node index.\allowbreak{}js}, which prints one line (source: demo-app/\allowbreak{}package.json, demo-app/\allowbreak{}index.js).

\section{The notify job}

\subsection{When it runs}

\begin{lstlisting}
if: >-
  failure() && needs.build-and-test.outputs.failure_kind == 'code' &&
  (github.event_name == 'push' || startsWith(github.head_ref, 'ai-fix-'))
permissions:
  contents: read
  pull-requests: write
  issues: write
env:
  MAX_ATTEMPTS: 3
\end{lstlisting}

(source: .github/\allowbreak{}workflows/\allowbreak{}demo-ci.yml)

The condition has three parts. \texttt{failure()} makes the job run only when an earlier job failed. The \texttt{failure\_\allowbreak{}kind} check stops it for infrastructure problems: the file's own comment says that \texttt{failure()} alone would also fire for a bound port, DNS or a full runner disk, and that the model would then invent a patch for a problem that is not in the source. The third part limits the event. A push to \texttt{main} always qualifies. A pull request qualifies only when its head branch starts with \texttt{ai-\allowbreak{}fix-\allowbreak{}}, which is the prefix the healer uses. A human's pull request therefore never wakes the healer.

The job asks for \texttt{pull-\allowbreak{}requests:\allowbreak{} write} and \texttt{issues:\allowbreak{} write}. The file does not say why, but its last step comments on a pull request, which is the likely reason.

\texttt{MAX\_\allowbreak{}ATTEMPTS} is 3. That is the total number of model attempts per heal PR: the first heal plus two retries (source: comment in .github/\allowbreak{}workflows/\allowbreak{}demo-ci.yml).

\subsection{Step 1: decide}

The step with id \texttt{decide} produces four outputs: \texttt{branch}, \texttt{commit}, \texttt{attempt} and \texttt{action}.

On a push the answer is fixed: the branch is \texttt{GITHUB\_\allowbreak{}REF\_\allowbreak{}NAME}, the commit is \texttt{GITHUB\_\allowbreak{}SHA}, \texttt{attempt} is 0 and \texttt{action} is \texttt{heal}.

On a pull request the step looks at the head commit of the PR, not at \texttt{github.\allowbreak{}sha}. For a pull request event, \texttt{github.\allowbreak{}sha} points at a merge commit that GitHub makes, and that commit does not carry the healer's message. The step fetches the head commit's message through the API and reads the trailer:

\begin{lstlisting}
attempt=$(printf '%s\n' "$msg" | sed -n 's/^Heal-Attempt: *\([0-9][0-9]*\) *$/\1/p' | tail -n 1)
\end{lstlisting}

(source: .github/\allowbreak{}workflows/\allowbreak{}demo-ci.yml, step "Decide")

The \texttt{Heal-\allowbreak{}Attempt:\allowbreak{} N} line is added by the healer to every commit it pushes, and the first heal is numbered 1 (source: tools/\allowbreak{}nodes/\allowbreak{}extract\_\allowbreak{}ai\_\allowbreak{}response.js). The step then picks one of three actions:

\begingroup\small
\begin{longtable}{>{\raggedright\arraybackslash}p{0.311\dimexpr\linewidth-6\tabcolsep\relax}>{\raggedright\arraybackslash}p{0.156\dimexpr\linewidth-6\tabcolsep\relax}>{\raggedright\arraybackslash}p{0.533\dimexpr\linewidth-6\tabcolsep\relax}}
\toprule
\textbf{Trailer found} & \textbf{Action} & \textbf{Meaning} \\
\midrule
\endhead
\bottomrule
\endfoot
\endlastfoot
None & \texttt{none} & A human pushed to the heal branch and owns it now. The step prints a notice and nothing else happens \\
N is at least \texttt{MAX\_\allowbreak{}ATTEMPTS} & \texttt{giveup} & The cap is reached \\
Otherwise & \texttt{heal} & Call the healer again with \texttt{attempt=N} \\
\bottomrule
\caption{Step 1: decide}\label{tab:05-1}\\
\end{longtable}
\endgroup

With the cap at 3, the commit numbering works out as follows. The first heal commit carries \texttt{Heal-\allowbreak{}Attempt:\allowbreak{} 1}. When CI fails on it, the notify job reads 1, which is below 3, and calls the healer, which pushes a commit marked 2. A second failure reads 2 and produces a commit marked 3. A third failure reads 3 and gives up. That is three model attempts in total.

The counter sits in the commit message, so re-running the workflow cannot reset it (source: comment in .github/\allowbreak{}workflows/\allowbreak{}demo-ci.yml).

\subsection{Step 2: trigger the webhook}

This step runs when \texttt{action} is \texttt{heal}. It posts to the address in the repository secret \texttt{N8N\_\allowbreak{}WEBHOOK\_\allowbreak{}URL}, with the shared secret in the header \texttt{X-\allowbreak{}Webhook-\allowbreak{}Secret}:

\begin{lstlisting}
curl -sS --fail --max-time 20 -X POST "${{ secrets.N8N_WEBHOOK_URL }}" \
  -H "Content-Type: application/json" \
  -H "X-Webhook-Secret: ${{ secrets.N8N_WEBHOOK_SECRET }}" \
  -d '{ "run_id": "...", "repo": "...", "branch": "...", "commit": "...",
        "actor": "...", "failure_kind": "...", "attempt": "...",
        "pr_number": "..." }' \
  || echo "::warning::n8n webhook unreachable - heal not triggered"
\end{lstlisting}

(shortened; each value is a \texttt{\$\{\{ \}\}} expression in the real file; source: .github/\allowbreak{}workflows/\allowbreak{}demo-ci.yml)

The fields and what n8n does with them are listed in Chapter 4. Three details stand out. The request has a 20 second limit (\texttt{-\allowbreak{}-\allowbreak{}max-\allowbreak{}time 20}). \texttt{-\allowbreak{}-\allowbreak{}fail} makes curl report an HTTP error as a failure. And the trailing \texttt{\textbar{}\textbar{} echo ":\allowbreak{}:\allowbreak{}warning:\allowbreak{}:\allowbreak{}.\allowbreak{}.\allowbreak{}.\allowbreak{}"} turns any failure of the call into a warning annotation, so the step itself still succeeds.

The warning exists because the webhook address is a tunnel that is only up while the healer is deployed. Without the fallback, a stopped healer would turn every broken commit red a second time, for a reason unrelated to the commit. The go-live notes record a test with the healer down: the build job failed and the notify job ended as success, with the warning shown in the annotations (source: docs/\allowbreak{}results/\allowbreak{}GOLIVE\_\allowbreak{}RESULTS.md, Round 2).

\subsection{Step 3: give up at the cap}

When \texttt{action} is \texttt{giveup}, the last step posts a comment on the pull request with \texttt{gh pr comment}. The text says how many attempts were made, that the pipeline has stopped and that the case needs a human (source: .github/\allowbreak{}workflows/\allowbreak{}demo-ci.yml). This is the only place where the cap is enforced. The n8n workflow does not check it (source: docs/\allowbreak{}architecture/\allowbreak{}n8n-workflow.md).

\section{The demo application}

The application is small on purpose. It stands in for a real project, and the pipeline would work the same way on one (source: demo-app/\allowbreak{}index.js, comment).

\begingroup\small
\begin{longtable}{>{\raggedright\arraybackslash}p{0.510\dimexpr\linewidth-4\tabcolsep\relax}>{\raggedright\arraybackslash}p{0.490\dimexpr\linewidth-4\tabcolsep\relax}}
\toprule
\textbf{File} & \textbf{Role} \\
\midrule
\endhead
\bottomrule
\endfoot
\endlastfoot
\texttt{demo-\allowbreak{}app/\allowbreak{}index.\allowbreak{}js} & Entry point. It prints one line and nothing else \\
\texttt{demo-\allowbreak{}app/\allowbreak{}src/\allowbreak{}app.\allowbreak{}js} & An Express app with three routes: \texttt{GET /\allowbreak{}health}, \texttt{GET /\allowbreak{}greet/\allowbreak{}:\allowbreak{}name} and \texttt{GET /\allowbreak{}sum?a=\&b=} \\
\texttt{demo-\allowbreak{}app/\allowbreak{}src/\allowbreak{}format.\allowbreak{}js} & Two helpers, \texttt{formatGreeting} and \texttt{formatSum}, used by the routes \\
\texttt{demo-\allowbreak{}app/\allowbreak{}test.\allowbreak{}js} & The smoke test: starts its own Express server on \texttt{PORT} (default 5000), requests \texttt{/\allowbreak{}} and expects status 200 \\
\texttt{demo-\allowbreak{}app/\allowbreak{}test/\allowbreak{}api.\allowbreak{}test.\allowbreak{}js} & The API test: loads \texttt{src/\allowbreak{}app.\allowbreak{}js}, listens on \texttt{API\_\allowbreak{}PORT} (default 5055) and checks four requests \\
\texttt{demo-\allowbreak{}app/\allowbreak{}package.\allowbreak{}json} & The one dependency is \texttt{express} \texttt{\textasciicircum{}4.\allowbreak{}19.\allowbreak{}2}. The \texttt{test} script is \texttt{node test.\allowbreak{}js \&\& node test/\allowbreak{}api.\allowbreak{}test.\allowbreak{}js} \\
\bottomrule
\caption{The demo application}\label{tab:05-2}\\
\end{longtable}
\endgroup

(source: the files listed)

The four API checks are \texttt{/\allowbreak{}health} returning \texttt{\{"status":\allowbreak{}"ok"\}}, \texttt{/\allowbreak{}greet/\allowbreak{}Ada} returning \texttt{Hello, Ada!}, \texttt{/\allowbreak{}sum?a=2\&b=3} returning a sum of 5, and \texttt{/\allowbreak{}sum?a=x\&b=3} returning status 400 with an error message (source: demo-app/\allowbreak{}test/\allowbreak{}api.test.js).

Both tests follow the same exit-code contract. A wrong status or body prints a line starting with \texttt{CODE\_\allowbreak{}FAILURE} and exits with 1. If the server cannot start, the \texttt{error} event on the server prints a line starting with \texttt{INFRA\_\allowbreak{}FAILURE} and exits with 75. For the smoke test, a failed \texttt{listen} such as \texttt{EADDRINUSE} goes to the infrastructure path, while a server that listens but does not answer itself counts as a code failure (source: demo-app/\allowbreak{}test.js, comments). We ran \texttt{npm test} with port 5000 already occupied and got exit code 75, which is the value the workflow turns into \texttt{infra}.

Because the script joins the two tests with \texttt{\&\&}, a failing smoke test stops the run and the API test never starts. The exit code of the whole \texttt{npm test} is the one from the failing command.

The log lines matter beyond the exit code. When the healer later downloads the job log, the first \texttt{\#\#[error]} line and the text above it are what the model reads, so a message that starts with \texttt{CODE\_\allowbreak{}FAILURE:\allowbreak{}} and names the route or the stack trace gives it something concrete (source: tools/\allowbreak{}nodes/\allowbreak{}build\_\allowbreak{}ai\_\allowbreak{}prompt.js, \texttt{trimCiLog}).

\section{Why the heal PR gets its own CI run}

In the first live round the workflow ran only on pushes to \texttt{main}. The fix pull requests were not tested at all, and the team checked them by running the file by hand (source: docs/\allowbreak{}results/\allowbreak{}GOLIVE\_\allowbreak{}RESULTS.md, "Still not true"). That left two holes. A fix that parsed but was wrong would sit in an open PR with no red mark. And there was no signal that could start a retry.

Adding the \texttt{pull\_\allowbreak{}request} trigger closed both holes. After the change, 5 of 5 heal PRs ran CI and all 5 were green (source: docs/\allowbreak{}results/\allowbreak{}GOLIVE\_\allowbreak{}RESULTS.md, Round 2). A red result on a heal PR is now the event that wakes the healer again, and the notify job's condition and the \texttt{Heal-\allowbreak{}Attempt} trailer keep that loop under control.

The \texttt{ai-\allowbreak{}fix-\allowbreak{}} prefix test has a second purpose, which is to stop the healer from feeding on its own mistakes without limit. A scratch pull request with a broken \texttt{test.\allowbreak{}js} was used to test this gate: the build job failed and the notify job was skipped, because the branch did not start with \texttt{ai-\allowbreak{}fix-\allowbreak{}} (source: docs/\allowbreak{}results/\allowbreak{}GOLIVE\_\allowbreak{}RESULTS.md, Round 2, run 35802293458). Separate scratch runs with trailers 3, 1 and none produced the give-up comment, a webhook call and nothing, in that order (source: docs/\allowbreak{}results/\allowbreak{}GOLIVE\_\allowbreak{}RESULTS.md, Round 3, "Gate tests without GPU").
